% ============================================================
%  Geometría Analítica · LM · Universidad de Sonora
%  Semana 1 — Presentación del curso y Unidad 1 (inicio)
%  Compilar en Overleaf con pdfLaTeX
% ============================================================
\documentclass[aspectratio=169,10pt]{beamer}
\usepackage[utf8]{inputenc}
\usepackage[T1]{fontenc}
\usepackage[spanish,es-nodecimaldot,es-noquoting]{babel}
\usepackage{lmodern}
\usepackage{amsmath,amssymb}
\usepackage{booktabs}
\usepackage{tikz}
\usetikzlibrary{arrows.meta,calc}
\usepackage{graphicx}
\graphicspath{{figuras/}}
% ---------- Tema y colores ----------
%  Misma estructura de plantilla que el curso de Estadística, con la
%  paleta propia de este curso: turquesa #0E7C8C (adoptado para todo
%  el material), con terracota y dorado como acentos.
\usetheme{default}
\setbeamertemplate{navigation symbols}{}
\definecolor{turquesaGA}{HTML}{0E7C8C}
\definecolor{turquesaClaro}{HTML}{E2F2F4}
\definecolor{grisT}{HTML}{5F5E5A}
\definecolor{negroP}{HTML}{2B2B2B}
\definecolor{terracota}{HTML}{C1502E}
\definecolor{dorado}{HTML}{B07C10}
\definecolor{verdeOk}{HTML}{1D9E75}
\definecolor{rojoAc}{HTML}{C0392B}
\setbeamercolor{structure}{fg=turquesaGA}
\setbeamercolor{frametitle}{fg=white,bg=turquesaGA}
\setbeamercolor{title}{fg=turquesaGA}
\setbeamerfont{frametitle}{size=\large,series=\bfseries}
\setbeamertemplate{frametitle}{%
  \nointerlineskip
  \begin{beamercolorbox}[wd=\paperwidth,ht=2.6ex,dp=1.2ex,leftskip=0.3cm]{frametitle}
    \usebeamerfont{frametitle}\insertframetitle
  \end{beamercolorbox}}
\setbeamertemplate{footline}{%
  \hfill\textcolor{grisT}{\tiny Geometría Analítica · LM · Unison \quad
  \insertframenumber/\inserttotalframenumber}\hspace{0.3cm}\vspace{0.15cm}}
\setbeamertemplate{itemize item}{\textcolor{turquesaGA}{\raisebox{0.15ex}{\tiny$\blacksquare$}}}
\setbeamertemplate{itemize subitem}{\textcolor{grisT}{\raisebox{0.15ex}{\tiny$\bullet$}}}
% ---------- Caja auxiliar para destacar ideas ----------
\newcommand{\destaca}[1]{%
  \par\vspace{3pt}\noindent
  \colorbox{turquesaClaro}{%
    \parbox{\dimexpr\linewidth-2\fboxsep\relax}{\vspace{3pt}#1\vspace{3pt}}}%
  \par\vspace{3pt}}
% ---------- Hojas de separación por sesión ----------
%  Cumplen dos funciones: al preparar la clase, ver de un golpe dónde terminó
%  cada día; al repasar, recordar en qué sesión se vio cada concepto.
%  Las sesiones no llevan número: se identifican por su título.
\newcommand{\hojaSesion}[2]{%
  \vfill
  \begin{beamercolorbox}[wd=\textwidth,sep=16pt,center]{frametitle}
    {\Large\bfseries #1}\\[11pt]
    \begin{minipage}{0.82\textwidth}\centering\small #2\end{minipage}
  \end{beamercolorbox}
  \vfill}
% ============================================================
\title{\LARGE\textbf{Geometría Analítica}}
\subtitle{Licenciatura en Matemáticas}
\author{Prof. Rosalía Gpe. Hernández Amador}
\institute{Departamento de Matemáticas · Universidad de Sonora}
\date{Semestre 2026-2}
\begin{document}
% ============================================================
\begin{frame}[plain]
  \titlepage
\end{frame}
% ============================================================
\section{Presentación del curso}
% ---------- hoja de sesión ----------
\begin{frame}[plain]
  \hojaSesion{Presentación del curso}{%
  Qué estudia la Geometría Analítica y cómo está organizado el curso.\\
  El temario y sus tres partes, la bibliografía y los criterios de evaluación.}
\end{frame}

% ============================================================
\begin{frame}{Datos del curso}
  \centering
  {\small
  \begin{tabular}{@{}ll@{}}
    \toprule
    Profesora & Rosalía Gpe.\ Hernández Amador \\
    Correo institucional & \texttt{rosalia.hernandez@unison.mx} \\
    Cubículo & Edificio 3K2, Módulo 9, Cubículo 1 \\
    \midrule
    Avisos, recordatorios y entregas & Teams \\
    Cuestionarios y prácticas & Moodle (Ramanujan) \\
%    Notas del curso & en Moodle, al cierre de cada semana \\
    \bottomrule
  \end{tabular}}
 
  \vspace{0.5cm}
  \raggedright
  \destaca{Toda la comunicación del curso es por medios institucionales;
  conviene revisar Teams periódicamente. 
  %Las notas de cada semana cubren exactamente lo visto en clase, sesión por sesión.
  }
\end{frame}

% ============================================================
\begin{frame}{¿De qué trata la Geometría Analítica?}
  Es el área de las matemáticas que estudia la relación entre las
  \textbf{figuras geométricas} y las \textbf{ecuaciones} que las
  describen:

  \vspace{0.2cm}
  \begin{center}
  \begin{tikzpicture}[>=Stealth]
    \node[fill=turquesaClaro,rounded corners,minimum height=1.05cm,
      minimum width=3.4cm] (F) at (0,0) {\textbf{figura}};
    \node[fill=turquesaClaro,rounded corners,minimum height=1.05cm,
      minimum width=3.4cm] (E) at (7.4,0) {\textbf{ecuación}};
    \draw[->,turquesaGA,very thick] ([yshift=7pt]F.east) -- ([yshift=7pt]E.west)
      node[midway,above,font=\small]{¿qué ecuación la describe?};
    \draw[->,terracota,very thick] ([yshift=-7pt]E.west) -- ([yshift=-7pt]F.east)
      node[midway,below,font=\small]{¿qué figura representa?};
  \end{tikzpicture}
  \end{center}
  
  \vspace{0.2cm}
  A lo largo del semestre aprenderemos a:

  \begin{itemize}
    \item \textbf{traducir} entre ambos lados, en las dos direcciones
      (que son distintas)
    \item \textbf{calcular} con álgebra lo que antes hacíamos con trazos
      (intersecciones, distancias, ángulos)
    \item \textbf{justificar} que la ecuación describe exactamente la
      figura, sabiendo por qué funcionan las fórmulas
  \end{itemize}

  \vspace{0.1cm}
  \destaca{La idea fundadora se debe a Descartes (1637):
  \textbf{asignar números a los puntos}, para tratar los problemas de
  figuras con las herramientas del álgebra. En este curso desarrollamos
  esa idea.}
\end{frame}

% ============================================================
\begin{frame}{Estructura del temario}
  Siete unidades, organizadas en tres partes:
  
%  \vspace{0.1cm}
  \begin{center}
  {\small
  \begin{tabular}{@{}cl@{}}
    \toprule
    \multicolumn{2}{@{}l}{\textcolor{turquesaGA}{\textbf{Parte I --- El plano}}} \\
    1 & Sistemas de coordenadas \\
    2 & La recta en el plano \\
    \midrule
    \multicolumn{2}{@{}l}{\textcolor{turquesaGA}{\textbf{Parte II --- El espacio}}} \\
    3 & Vectores en el espacio \\
    4 & Planos y rectas en el espacio \\
    \midrule
    \multicolumn{2}{@{}l}{\textcolor{turquesaGA}{\textbf{Parte III --- Curvas y superficies de segundo grado}}} \\
    5 & Cónicas \\
    6 & Superficies cuádricas \\
    7 & Otros sistemas de coordenadas \\
    \bottomrule
  \end{tabular}}
  \end{center}
  
 % \vspace{0.1cm}
  En la parte I desarrollamos la teoría con los objetos más simples; la
  parte II la lleva al espacio, donde el lenguaje de \textit{vectores}
  muestra su ventaja; la parte III la aplica a las curvas y
  superficies de segundo grado.\\

  \vspace{0.15cm}
  {\footnotesize El temario completo se encuentra en la página de la carrera,
  \href{https://lic.mat.uson.mx/}{\texttt{lic.mat.uson.mx}}.}
\end{frame}

% ============================================================
\begin{frame}{Objetivos del curso}
  \begin{center}
  {\small
  \begin{tabular}{@{}p{6.4cm}p{6.6cm}@{}}
    \toprule
    \textcolor{rojoAc}{\textbf{No se trata de\ldots}} &
    \textcolor{verdeOk}{\textbf{sino de\ldots}} \\
    \midrule
    Memorizar un formulario de rectas y cónicas &
    Saber \textbf{traducir} en ambas direcciones:
    figura $\leftrightarrow$ ecuación \\[8pt]
    Aplicar fórmulas sin saber de dónde salen &
    \textbf{Justificar} por qué funciona cada fórmula \\[8pt]
    Confundir un dibujo bien hecho con una demostración &
    \textbf{Verificar} los resultados por más de una vía \\
    \bottomrule
  \end{tabular}}
  \end{center}
  
  \vspace{0.3cm}
%  \destaca{Estos tres verbos (\textbf{traducir}, \textbf{justificar},
%  \textbf{verificar}) reaparecerán en los criterios de los exámenes.}
\end{frame}

% ============================================================
\begin{frame}{Bibliografía}
  Existe mucha literatura de Geometría Analítica, y toda buena fuente
  sirve para consultar. Nos apoyamos en estos materiales:

  \vspace{0.3cm}
  \begin{center}
  {\small
  \begin{tabular}{@{}p{2.5cm}p{7.6cm}@{}}
    \toprule
    \textbf{Texto base} & Ramírez Galarza, \textit{Geometría analítica.
      Una introducción a la geometría}, Prensas de Ciencias, UNAM \\[6pt]
    \textbf{Problemas} & Lehmann, \textit{Geometría analítica}, Limusa.
      Es un clásico; sus series de ejercicios son abundantes y avanzan
      de lo sencillo a lo difícil \\[6pt]
    \textbf{Adicional} & Bracho, \textit{Introducción analítica
      a las geometrías}, FCE \\[6pt]
    \textbf{Libro digital}\newline
      {\footnotesize\itshape (en construcción)} & Cubre el material en el orden visto en
      clase e incluye figuras interactivas.\newline
      {\footnotesize\href{https://rosaliahdez.github.io/notas-geometria-analitica/}{\texttt{rosaliahdez.github.io/notas-geometria-analitica}}} \\
    \bottomrule
  \end{tabular}}
  \end{center}

\end{frame}

% ============================================================
%  Evaluación acordada con el grupo: exámenes parciales 100 %;
%  todo lo demás abona puntos extra.
\begin{frame}{Criterios de evaluación del curso}
  \centering
  {\small
  \begin{tabular}{@{}p{5.4cm}c p{6cm}@{}}
    \toprule
    \textbf{Instrumento} & \textbf{Valor} & \textbf{Notas} \\
    \midrule
    Exámenes parciales & 100\,\% & En papel; se toma en cuenta el procedimiento \\
    Prácticas de laboratorio (GeoGebra) & pase & Se entregan en Moodle; pase para el parcial \\
    Cuestionarios y tareas & pase & En Moodle y en papel; pase para el parcial \\
    \midrule
    \textbf{Total} & \textbf{100\,\%} &  \\
    \bottomrule
  \end{tabular}}

  \vspace{0.5cm}
  \raggedright
  \destaca{En los exámenes no se evalúa la memorización. Se evalúa que
  sepan \textbf{traducir} entre figura y ecuación, \textbf{calcular}
  con orden, \textbf{justificar} los pasos, y \textbf{verificar} el
  resultado por una segunda vía cuando sea posible.}

\vspace{0.3cm}
  {\footnotesize Para tener derecho a evaluación ordinaria se requiere al
  menos el \textbf{75\,\% de asistencia}, conforme al Artículo 70 del
  Reglamento Escolar de la Universidad de Sonora.}
\end{frame}

% ============================================================
%\begin{frame}{Dos acuerdos de trabajo}
%  \textbf{1. Predicción antes de observación.} Cada vez que proyectemos
%  una figura (estática o dinámica) primero diremos qué esperamos ver,
%  y después la veremos. Una figura que confirma lo que ya pensábamos
%  enseña poco; una que nos
%  \textcolor{terracota}{\textbf{contradice}} enseña muchísimo.
%  \vspace{0.4cm}

%  \textbf{2. Las dos direcciones.} De la ecuación a la figura
%  \textbf{y} de la figura a la ecuación. Son habilidades distintas
%  (dominar una no garantiza la otra), así que las practicaremos y las
%  evaluaremos por separado.
%  \vspace{0.4cm}
%  \destaca{Estos dos acuerdos aparecerán cada semana en los applets, en
%  los talleres y en los criterios de los exámenes.}
%\end{frame}
% ============================================================


\section{Puntos y distancia}
% ---------- Portada de la Unidad 1 ----------
%  Página ligera de apertura de unidad. La hoja verde de sesión viene
%  justo después con su resumen; si se prefiere una sola página, basta
%  borrar esta o aquella.
\begin{frame}[plain]
  \begin{columns}[c]
    \begin{column}{0.42\textwidth}
      \centering
      \begin{tikzpicture}[scale=0.6,>=Stealth]
        \draw[very thin,color=gray!30] (-0.6,-0.6) grid (5.4,3.6);
        \draw[->] (-0.6,0) -- (5.6,0) node[below,font=\small]{$x$};
        \draw[->] (0,-0.6) -- (0,3.8) node[left,font=\small]{$y$};
        \draw[turquesaGA,very thick] (1,1) -- (4.5,3);
        \draw[dashed,terracota,thick] (1,1) -- (4.5,1) -- (4.5,3);
        \fill[turquesaGA] (1,1) circle(3pt);
        \fill[turquesaGA] (4.5,3) circle(3pt);
      \end{tikzpicture}
    \end{column}
    \begin{column}{0.52\textwidth}
      \raggedright
      {\small\color{grisT} SEMANA 1}\\[8pt]
      {\color{dorado}\rule{0.55\linewidth}{1.6pt}}\\[12pt]
      {\LARGE\bfseries\color{turquesaGA} Unidad 1}\\[8pt]
      {\large Sistemas de coordenadas}\\[14pt]
      {\small\color{grisT}\itshape Aprender a traducir entre la figura
      y la ecuación}
    \end{column}
  \end{columns}
\end{frame}

% ---------- hoja de sesión ----------
\begin{frame}[plain]
  \hojaSesion{La geometría analítica.\\ Puntos y distancia en $\mathbb{R}^2$}{%
  Las geometrías del plan de estudios y qué distingue a la analítica.
  El plano cartesiano: ejes, cuadrantes y el papel de los signos.
  La fórmula de la distancia y su origen geométrico en el teorema de
  Pitágoras.}
\end{frame}
\begin{frame}{Las geometrías del plan de estudios}
  \begin{itemize}
    \item La \textbf{euclidiana}, que ya conocen del bachillerato:
      figuras, congruencia, regla y compás. Se estudia razonando
      directamente sobre las figuras.

      \vspace{0.2cm}
    \item Más adelante en la carrera: proyectiva, diferencial, \dots

      \vspace{0.2cm}
    \item La \textbf{analítica} (este curso) estudia la relación

      \vspace{2mm}
      \begin{center}
        {\Large\color{turquesaGA} figura $\longleftrightarrow$ ecuación}
      \end{center}
  \end{itemize}
  
  \vspace{0.4cm}
  \destaca{Antes de Descartes, que dos rectas se cortan se comprobaba
  trazándolas. Después de Descartes, lo comprobamos \textbf{resolviendo
  un sistema de ecuaciones}, algo que podemos hacer de forma algebraica,
  sistemática y sin depender de un dibujo bien hecho. Ese cambio de
  método es el tema de este curso.}
\end{frame}

% ============================================================
\begin{frame}{El plano cartesiano}
  \begin{columns}[T]
    \begin{column}{0.52\textwidth}
      \begin{itemize}
        \item Dos rectas perpendiculares: los \textbf{ejes} $X$ e $Y$
        \item Su intersección: el \textbf{origen} $O=(0,0)$
        \item Cada punto corresponde a una \textbf{pareja ordenada}
          $(x,y)$; $x$ es la \textbf{abscisa} y $y$ la \textbf{ordenada}
        \item Los ejes determinan cuatro \textbf{cuadrantes}, numerados
          en sentido contrario a las manecillas del reloj
      \end{itemize}
      
      \vspace{0.3cm}
      \destaca{\small Que la \textbf{pareja sea ordenada} significa que
      el orden importa: $(3,5)\neq(5,3)$.}
    \end{column}
    \begin{column}{0.44\textwidth}
      \begin{center}
      \begin{tikzpicture}[scale=.62,>=Stealth]
        \fill[turquesaClaro] (0,0) rectangle (3.2,2.6);
        \draw[very thin,color=gray!30] (-3.2,-2.6) grid (3.2,2.6);
        \draw[->,thick] (-3.2,0) -- (3.4,0) node[below]{$X$};
        \draw[->,thick] (0,-2.6) -- (0,2.8) node[left]{$Y$};
        \node[color=turquesaGA] at (0.9,2.1) {\textbf{I}};
        \node[color=grisT] at (-1.6,1.4) {\textbf{II}};
        \node[color=grisT] at (-1.6,-1.4) {\textbf{III}};
        \node[color=grisT] at (1.6,-1.4) {\textbf{IV}};
        \draw[dashed,terracota,thick] (2,1.5) -- (2,0);
        \draw[dashed,terracota,thick] (2,1.5) -- (0,1.5);
        \fill[terracota] (2,1.5) circle(2.6pt) node[above right]{\small $(2,1.5)$};
        \node[above left,font=\small] at (0,0) {$O$};
      \end{tikzpicture}
      \end{center}
    \end{column}
  \end{columns}
\end{frame}
% ============================================================
\begin{frame}{Signos y cuadrantes}
  \begin{columns}[T]
    \begin{column}{0.56\textwidth}
      \centering
      {\small
      \begin{tabular}{@{}lccl@{}}
        \toprule
        \textbf{Región} & \textbf{signo de $x$} & \textbf{signo de $y$} & \\
        \midrule
        Cuadrante I & $+$ & $+$ & \\
        Cuadrante II & $-$ & $+$ & \\
        Cuadrante III & $-$ & $-$ & \\
        Cuadrante IV & $+$ & $-$ & \\
        \midrule
        Eje $X$ & cualquiera & $\mathbf{0}$ & ordenada nula \\
        Eje $Y$ & $\mathbf{0}$ & cualquiera & abscisa nula \\
        \bottomrule
      \end{tabular}}
    \end{column}
    \begin{column}{0.40\textwidth}
      \begin{center}
      \begin{tikzpicture}[scale=.62,>=Stealth]
        %\fill[turquesaClaro] (0,0) rectangle (3.2,2.6);
        \draw[very thin,color=gray!30] (-3.2,-2.6) grid (3.2,2.6);
        \draw[->,thick] (-3.2,0) -- (3.4,0) node[below]{$X$};
        \draw[->,thick] (0,-2.6) -- (0,2.8) node[left]{$Y$};
        \node[color=grisT] at (1.6,1.4) {\textbf{I}};
        \node[color=grisT] at (-1.6,1.4) {\textbf{II}};
        \node[color=grisT] at (-1.6,-1.4) {\textbf{III}};
        \node[color=grisT] at (1.6,-1.4) {\textbf{IV}};
        \node[below left,font=\small] at (0,0) {$O$};
      \end{tikzpicture}
      \end{center}
    \end{column}
  \end{columns}

  \vspace{0.4cm}
  \raggedright
  \destaca{La tabla también funciona \textbf{al revés}: conocer los signos
  de un punto (o de expresiones como $xy$) restringe dónde puede estar.
  Por ejemplo, si $xy<0$, los signos son opuestos y el punto está en el
  cuadrante II o en el IV. %\Este razonamiento apareceré en el cuestionario de la semana.
  }
\end{frame}
% ============================================================
\begin{frame}{Distancia del origen a un punto}
  \begin{columns}[T]
    \begin{column}{0.52\textwidth}
      Empezamos por el caso más sencillo: hallar la distancia de un
      punto $P=(x_1,y_1)$ del primer cuadrante al origen.

      \vspace{0.15cm}
      El segmento $OP$ es la \textbf{hipotenusa}
      de un triángulo rectángulo cuyos catetos son las coordenadas
      mismas de $P$.

%      \vspace{0.15cm}
      \begin{itemize}
        \item El cateto horizontal está sobre el eje $X$ y mide $x_1$.
        \item El cateto vertical mide $y_1$.
        \item El ángulo recto está en $(x_1,0)$.
      \end{itemize}

 %     \vspace{0.15cm}
      Por el teorema de Pitágoras,
      \vspace{-0.2cm}
      \[
        d(O,P)=\sqrt{x_1^{\,2}+y_1^{\,2}} .
      \]
    \end{column}
    \begin{column}{0.44\textwidth}
      \begin{center}
      \begin{tikzpicture}[scale=.72,>=Stealth]
        \draw[very thin,color=gray!30] (-0.4,-0.4) grid (4.4,3.2);
        \draw[->] (-0.4,0) -- (4.6,0) node[below]{$x$};
        \draw[->] (0,-0.4) -- (0,3.4) node[left]{$y$};
        \coordinate (O) at (0,0);
        \coordinate (P) at (3.4,2.4);
        \coordinate (R) at (3.4,0);
        \fill[turquesaGA!10] (O) -- (R) -- (P) -- cycle;
        \draw[dashed,terracota,thick] (O) -- (R)
          node[midway,below]{\small $x_1$};
        \draw[dashed,terracota,thick] (R) -- (P)
          node[midway,right]{\small $y_1$};
        \draw[turquesaGA,very thick] (O) -- (P)
          node[midway,above,sloped]{\small $d(O,P)$};
        \draw (3.1,0) -- (3.1,0.3) -- (3.4,0.3);
        \fill[turquesaGA] (P) circle(2.2pt) node[above]{$P$};
        \node[below left,font=\small] at (O) {$O$};
      \end{tikzpicture}
      \end{center}
    \end{column}
  \end{columns}

  \vspace{0.05cm}
  \destaca{\textbf{Ejemplo 1.} $d\bigl(O,(3,4)\bigr)=\sqrt{9+16}
  =\sqrt{25}=5$. En los otros cuadrantes el razonamiento es igual, con
  catetos $|x_1|$ y $|y_1|$; los valores absolutos desaparecen al
  elevar al cuadrado.}
\end{frame}
% ============================================================
\begin{frame}{El mismo triángulo, fuera del origen}
  \begin{columns}[T]
    \begin{column}{0.50\textwidth}
      Ahora ninguno de los dos puntos está en el origen. Tomamos
      $P=(x_1,y_1)$ y $Q=(x_2,y_2)$, y trasladamos el triángulo. Ahora el vértice del ángulo recto está en $(x_2,y_1)$.

      \vspace{0.15cm}
      Los catetos ya no son coordenadas sino \textbf{diferencias de
      coordenadas}, y miden $|x_2-x_1|$ y $|y_2-y_1|$.
      \[
        d(P,Q)=\sqrt{(x_2-x_1)^2+(y_2-y_1)^2}
      \]

      \vspace{0.05cm}
      Obsérvese que si $P$ es el origen, entonces $x_1=y_1=0$ y recuperamos la
      fórmula anterior.
    \end{column}
    \begin{column}{0.46\textwidth}
      \begin{center}
      \begin{tikzpicture}[scale=.63,>=Stealth]
        \draw[very thin,color=gray!30] (-0.4,-0.4) grid (5.6,3.4);
        \draw[->] (-0.4,0) -- (5.8,0) node[below]{$x$};
        \draw[->] (0,-0.4) -- (0,3.6) node[left]{$y$};
        \coordinate (P) at (1,1);
        \coordinate (Q) at (4.5,3);
        \coordinate (R) at (4.5,1);
        \fill[turquesaGA!10] (P) -- (R) -- (Q) -- cycle;
        \node[below,font=\scriptsize] at (1,0) {$x_1$};
        \node[below,font=\scriptsize] at (4.5,0) {$x_2$};
        \node[left,font=\scriptsize] at (0,1) {$y_1$};
        \node[left,font=\scriptsize] at (0,3) {$y_2$};
        \draw[turquesaGA,very thick] (P) -- (Q)
          node[midway,above,sloped]{\small $d(P,Q)$};
        \draw[dashed,terracota,thick] (P) -- (R)
          node[midway,below]{\small $|x_2-x_1|$};
        \draw[dashed,terracota,thick] (R) -- (Q)
          node[midway,right]{\small $|y_2-y_1|$};
        \draw (4.2,1) -- (4.2,1.3) -- (4.5,1.3);
        \fill[turquesaGA] (P) circle(2.2pt) node[above left]{$P$};
        \fill[turquesaGA] (Q) circle(2.2pt) node[above]{$Q$};
      \end{tikzpicture}
      \end{center}
    \end{column}
  \end{columns}

  \vspace{0.15cm}
  \destaca{\small La fórmula \textbf{no es una definición}, es el
  teorema de Pitágoras midiendo la hipotenusa $PQ$. Tiene esta forma
  porque los ejes son perpendiculares y las unidades iguales en ambos.
  Con ejes oblicuos, la distancia se calcularía de otro modo.}
\end{frame}
% ============================================================
\begin{frame}{Ejemplos}
  \textbf{Ejemplo 2.} $d\bigl((1,-2),\,(4,2)\bigr)
    =\sqrt{3^2+4^2}=\sqrt{25}=5$.

  \vspace{0.4cm}

  \textbf{Ejemplo 3.} $d\bigl((-3,1),\,(2,-1)\bigr)
    =\sqrt{5^2+(-2)^2}=\sqrt{29}$.

  \vspace{0.4cm}

  \textbf{Ejemplo 4.} Con $A=(1,1)$, $B=(5,3)$ y $C=(3,7)$ obtenemos
  $d(A,B)=\sqrt{20}$, $d(B,C)=\sqrt{20}$ y $d(A,C)=\sqrt{40}$. El
  triángulo $ABC$ es isósceles, y como
  $d(A,B)^2+d(B,C)^2=d(A,C)^2$, el ángulo recto está en $B$.

  \vspace{0.4cm}
  \destaca{\textbf{Convención del curso: las raíces se dejan indicadas.}
  $\sqrt{29}$ es una respuesta exacta; $5.385\ldots$ es una aproximación.
  Preferimos las respuestas exactas, y así se pedirán en las
  evaluaciones.}
\end{frame}

% ============================================================
\section{Propiedades de la distancia y lugares geométricos}
% ---------- hoja de sesión ----------
\begin{frame}[plain]
  \hojaSesion{Propiedades de la distancia.\\ Lugares geométricos}{%
  Las propiedades que cumple toda distancia y por qué resultan
  creíbles geométricamente antes de cualquier cálculo. El concepto de
  lugar geométrico (central en todo el curso) y los primeros
  ejemplos: rectas paralelas a los ejes.}
\end{frame}

% ============================================================
\begin{frame}{Propiedades de la distancia}
  Para cualesquiera puntos $P$, $Q$, $R$:

  \vspace{0.2cm}
  \begin{enumerate}
    \item $d(P,Q)\ \ge\ 0$, \ y \ $d(P,Q)=0 \iff P=Q$
    \item $d(P,Q)=d(Q,P)$ \hfill (\textit{simetría})
    \item $d(P,Q)\ \le\ d(P,R)+d(R,Q)$ \hfill (\textit{desigualdad del triángulo})
  \end{enumerate}

  \vspace{0.4cm}
  \destaca{\textbf{¿Por qué son creíbles antes de calcular nada?}
  \begin{enumerate}
  \item Una longitud no puede ser negativa, y la única forma de recorrer
  distancia cero es no moverse.
  \item El camino de ida mide lo mismo que el
  de vuelta. 
  \item Y pasar por un punto intermedio $R$ no puede
  \textbf{acortar} el trayecto de $P$ a $Q$: en el mejor de los casos lo
  iguala, exactamente cuando $R$ está en el segmento $PQ$.
  \end{enumerate}}
\end{frame}
% ============================================================
\begin{frame}{Ahora, con la fórmula}
  Lo que resulta plausible geométricamente, hay que saberlo verificar:

  \vspace{0.2cm}
  \begin{itemize}
    \item[1.] Una raíz de suma de cuadrados nunca es negativa, y vale
      $0$ únicamente cuando \textbf{ambos} cuadrados son $0$: es decir,
      cuando $x_1=x_2$ y $y_1=y_2$.
    \item[2.] $(x_2-x_1)^2=(x_1-x_2)^2$: elevar al cuadrado elimina el
      signo, así que intercambiar $P$ con $Q$ no cambia el valor.
    \item[3.] La desigualdad del triángulo es más trabajosa: requiere
      una herramienta que construiremos más adelante (el producto punto).
      Queda \textbf{enunciada, no demostrada}.
  \end{itemize}

  \vspace{0.4cm}
  \destaca{Registrar qué está demostrado y qué está pendiente
  \textbf{es parte del trabajo matemático}. Dejaremos anotado el
  pendiente y lo demostraremos más adelante.}
\end{frame}

% ============================================================
\begin{frame}{Lugar geométrico}
  \destaca{\textbf{Definición.} Un \textbf{lugar geométrico} es el
  conjunto de \textbf{todos} los puntos que cumplen una condición dada,
  y \textbf{solamente} ellos.}

  \vspace{0.4cm}
  Las dos palabras esenciales son \textit{todos} y \textit{solamente}.
  Afirmar que una ecuación describe un lugar geométrico es afirmar
  \textbf{dos} cosas a la vez:

  \vspace{0.2cm}
  \begin{enumerate}
    \item todo punto que cumple la condición satisface la ecuación;
    \item todo punto que satisface la ecuación cumple la condición.
  \end{enumerate}

  \vspace{0.3cm}
  {\footnotesize Omitir la segunda mitad es un error clásico: por
  ejemplo, elevar al cuadrado puede introducir soluciones que no
  cumplían la condición original. Debemos tener cuidado con esto a lo largo del semestre.}
\end{frame}

% ============================================================
\begin{frame}{Rectas paralelas a los ejes}
  \begin{columns}[T]
    \begin{column}{0.52\textwidth}
      \begin{itemize}
        \item El lugar geométrico en el plano de los puntos con abscisa $3$:
          \[
            x=3 \quad\text{\small (recta vertical, paralela al eje $Y$)}
          \]
        \item El lugar geométrico de los puntos con ordenada $-2$:
          \[
            y=-2 \quad\text{\small (recta horizontal, paralela al eje $X$)}
          \]
        \item Los propios ejes son lugares geométricos:
          \[
            \text{eje } X:\; y=0 \qquad \text{eje } Y:\; x=0
          \]
      \end{itemize}
    \end{column}
    \begin{column}{0.44\textwidth}
      \begin{center}
      \begin{tikzpicture}[scale=.6,>=Stealth]
        \draw[very thin,color=gray!30] (-2.6,-3.1) grid (4.1,2.6);
        \draw[->,thick] (-2.6,0) -- (4.3,0) node[below]{$X$};
        \draw[->,thick] (0,-3.1) -- (0,2.8) node[left]{$Y$};
        \draw[turquesaGA,very thick] (3,-3.1) -- (3,2.6)
          node[above,font=\small]{$x=3$};
        \draw[terracota,very thick] (-2.6,-2) -- (4.1,-2)
          node[above left,font=\small]{$y=-2$};
        \fill[dorado] (3,-2) circle(2.6pt);
      \end{tikzpicture}
      \end{center}
    \end{column}
  \end{columns}

  \vspace{0.15cm}
  \destaca{\small La descripción verbal es tan clara que la ecuación
  se obtiene con facilidad: la variable que \textbf{no aparece} en la
  ecuación es la que queda \textbf{libre}.

  \vspace{0.1cm}
  Nótese además el punto dorado de la figura, en la intersección de
  las rectas de colores: ahí se cumplen \textbf{las dos condiciones} a
  la vez. Una condición deja una recta; dos condiciones dejan un solo
  punto.}
\end{frame}

% ============================================================
\begin{frame}{Ejercicios}
  \textbf{De la condición a la ecuación:}
  \begin{enumerate}
    \item Los puntos que están a distancia $2$ del eje $X$, por encima
      de él.
    \item La recta paralela al eje $Y$ que pasa por $(4,-1)$.
  \end{enumerate}

  \vspace{0.4cm}
  \textbf{De la ecuación a la figura:}
  \begin{enumerate}
    \setcounter{enumi}{2}
    \item ¿Qué describe $y=0$? ¿Y el par de condiciones simultáneas
      $x=3$, $y=-2$?
  \end{enumerate}
  
  \vspace{0.4cm}
  \destaca{\small \textit{Soluciones.} 1.~$y=2$.\quad
  2.~$x=4$; la ordenada $-1$ no interviene, porque la condición
  ``paralela al eje $Y$'' solo fija la abscisa.\quad
  3.~El eje $X$; y el punto $(3,-2)$: dos condiciones reducen el lugar
  a un solo punto.}
\end{frame}

% ============================================================
\section{El espacio tridimensional}
% ---------- hoja de sesión ----------
\begin{frame}[plain]
  \hojaSesion{El espacio tridimensional $\mathbb{R}^3$}{%
  El espacio $\mathbb{R}^3$ construido en paralelo con $\mathbb{R}^2$,
  concepto por concepto: tres ejes perpendiculares y la convención
  para dibujarlos, los planos coordenados, los ocho octantes, y cómo
  dibujar un punto sin perder la perspectiva. Primeras preguntas de
  lugares geométricos en el espacio.}
\end{frame}


% ============================================================
\begin{frame}{De $\mathbb{R}^2$ a $\mathbb{R}^3$}
  Construimos el espacio apoyándonos en lo que ya sabemos del plano.
  Cada concepto tiene su análogo:
  
  \vspace{0.2cm}
  \begin{center}
  {\small
  \begin{tabular}{@{}ll@{}}
    \toprule
    \textbf{En el plano} & \textbf{En el espacio} \\
    \midrule
    pareja ordenada $(x,y)$ & terna ordenada $(x,y,z)$ \\
    dos ejes perpendiculares & tres ejes perpendiculares \\
    los ejes dividen el plano & los \textbf{planos coordenados}
      dividen el espacio \\
    cuatro cuadrantes & ocho \textbf{octantes} \\
    \bottomrule
  \end{tabular}}
  \end{center}
  
  \vspace{0.3cm}
  \destaca{Lo que construimos en $\mathbb{R}^2$ se
  traduce a $\mathbb{R}^3$ con los mismos argumentos. Cuando algo
  \textbf{no} se traduzca, lo diremos explícitamente.}
\end{frame}

% ============================================================
\begin{frame}{Convención para dibujar los ejes}
  \begin{columns}[T]
    \begin{column}{0.52\textwidth}
      Siempre dibujaremos los ejes así (sus partes positivas):

      \vspace{0.2cm}
      \begin{itemize}
        \item el eje $x$ apunta \textbf{hacia nosotros}
        \item el eje $y$ apunta \textbf{a la derecha}
        \item el eje $z$ apunta \textbf{hacia arriba}
      \end{itemize}
      
      \vspace{0.3cm}
      \destaca{\small Para ubicarse, sirve mirar la \textbf{esquina
      del salón}: las dos paredes y el piso se cortan como los tres
      planos coordenados, y la esquina es el origen.}
    \end{column}
    \begin{column}{0.44\textwidth}
      \begin{center}
      \begin{tikzpicture}[scale=1.05,>=Stealth]
        \draw[->,thick] (0,0) -- (-1.6,-1.28) node[below left]{$x$};
        \draw[->,thick] (0,0) -- (2.6,0) node[right]{$y$};
        \draw[->,thick] (0,0) -- (0,2.4) node[above]{$z$};
        \draw[gray!60] (0,0) -- (1.1,0.88);
        \draw[gray!60] (0,0) -- (-1.4,0);
        \draw[gray!60] (0,0) -- (0,-1.1);
        \node[below left,font=\small] at (0,0) {$O$};
      \end{tikzpicture}
      \end{center}
      
     \vspace{-0.4cm}
      {\footnotesize En gris, las partes negativas de cada eje.}
    \end{column}
  \end{columns}
  
  \vspace{0.4cm}
  {\footnotesize Se sugiere visitar el sitio de \textbf{GeoGebra 3D},
  \url{https://www.geogebra.org/3d}, para mover la perspectiva y ver
  los ejes desde cualquier lado.}
\end{frame}
% ============================================================
\begin{frame}{Planos coordenados}
  \begin{columns}[T]
    \begin{column}{0.50\textwidth}
      Cada pareja de ejes genera un \textbf{plano coordenado}: el
      plano $xy$, el plano $xz$ y el plano $yz$.

      \vspace{0.2cm}
      En la figura se ven las tres ``hojas'' completas (partes
      positivas y negativas) intersectándose en el origen.

      \vspace{0.3cm}
      \destaca{\small En el plano, dos rectas (los
      ejes) producen cuatro regiones; en el espacio, tres planos (los coordenados)
      producen ocho. La tabla de signos del plano tiene ahora una
      tercera columna.}
    \end{column}
    \begin{column}{0.46\textwidth}
      \begin{center}
      \begin{tikzpicture}[scale=0.72,>=Stealth]
        % hoja yz (al fondo)
        \fill[dorado!20,opacity=0.5] (-1.8,-1.3) rectangle (1.8,1.3);
        % hoja xz
        \fill[terracota!15,opacity=0.5]
          (-0.77,0.68) -- (0.77,1.92) -- (0.77,-0.68) -- (-0.77,-1.92) -- cycle;
        % hoja xy (al frente)
        \fill[turquesaClaro,opacity=0.5]
          (1.03,-0.62) -- (2.57,0.62) -- (-1.03,0.62) -- (-2.57,-0.62) -- cycle;
        % ejes: positivos en negro, negativos en gris
        \draw[->,thick] (0,0) -- (-1.15,-0.92) node[below left]{$x$};
        \draw[gray!60] (0,0) -- (0.99,0.79);
        \draw[->,thick] (0,0) -- (2.75,0) node[right]{$y$};
        \draw[gray!60] (0,0) -- (-2.0,0);
        \draw[->,thick] (0,0) -- (0,2.1) node[above]{$z$};
        \draw[gray!60] (0,0) -- (0,-1.6);
        \node[font=\scriptsize,color=grisT] at (1.75,-0.30) {plano $xy$};
        \node[font=\scriptsize,color=grisT] at (1.35,1.05) {plano $yz$};
        \node[font=\scriptsize,color=grisT] at (0.45,1.55) {plano $xz$};
      \end{tikzpicture}
      \end{center}
    \end{column}
  \end{columns}
\end{frame}
% ============================================================
\begin{frame}{Los ocho octantes}
  \begin{columns}[T]
    \begin{column}{0.50\textwidth}
      Los tres planos coordenados dividen el espacio en \textbf{ocho
      octantes}, como ocho cajas apiladas: cuatro arriba del plano
      $xy$ y cuatro abajo.

      \vspace{0.2cm}
      \begin{itemize}
        \item Octantes I a IV: arriba ($z>0$), empezando donde
          $x>0$, $y>0$ y avanzando contra las manecillas del reloj
        \item Octantes V a VIII: abajo ($z<0$), cada uno debajo del
          correspondiente
      \end{itemize}

      \vspace{0.3cm}
      \destaca{\small Los puntos de los ejes y de los planos
      coordenados no están en ningún octante, igual que en el plano
      los puntos de los ejes no están en ningún cuadrante.}
    \end{column}
    \begin{column}{0.46\textwidth}
      \begin{center}
      \begin{tikzpicture}[scale=0.82,>=Stealth]
        % aristas del cubo exterior
        \draw[gray!55] (0.55,0.66) -- (1.65,1.54) -- (-0.55,1.54) -- (-1.65,0.66) -- cycle;
        \draw[gray!55] (0.55,-1.54) -- (1.65,-0.66) -- (-0.55,-0.66) -- (-1.65,-1.54) -- cycle;
        \draw[gray!55] (0.55,0.66) -- (0.55,-1.54);
        \draw[gray!55] (1.65,1.54) -- (1.65,-0.66);
        \draw[gray!55] (-0.55,1.54) -- (-0.55,-0.66);
        \draw[gray!55] (-1.65,0.66) -- (-1.65,-1.54);
        % cortes de los planos coordenados
        \draw[turquesaGA!60,thin] (0.55,-0.44) -- (1.65,0.44) -- (-0.55,0.44) -- (-1.65,-0.44) -- cycle;
        \draw[turquesaGA!60,thin] (-0.55,0.66) -- (0.55,1.54) -- (0.55,-0.66) -- (-0.55,-1.54) -- cycle;
        \draw[turquesaGA!60,thin] (1.1,1.1) -- (-1.1,1.1) -- (-1.1,-1.1) -- (1.1,-1.1) -- cycle;
        % ejes
        \draw[->,thick] (0,0) -- (-0.99,-0.79) node[below left,font=\small]{$x$};
        \draw[->,thick] (0,0) -- (1.98,0) node[right,font=\small]{$y$};
        \draw[->,thick] (0,0) -- (0,1.87) node[above,font=\small]{$z$};
        % numeros: al frente en color, atras en gris
        \node[font=\small\bfseries,color=turquesaGA] at (0.275,0.33) {I};
        \node[font=\small,color=gray!70] at (0.825,0.77) {II};
        \node[font=\small,color=gray!70] at (-0.275,0.77) {III};
        \node[font=\small\bfseries,color=turquesaGA] at (-0.825,0.33) {IV};
        \node[font=\small\bfseries,color=turquesaGA] at (0.275,-0.77) {V};
        \node[font=\small,color=gray!70] at (0.825,-0.33) {VI};
        \node[font=\small,color=gray!70] at (-0.275,-0.33) {VII};
        \node[font=\small\bfseries,color=turquesaGA] at (-0.825,-0.77) {VIII};
      \end{tikzpicture}
      \end{center}

      {\footnotesize En turquesa, los números de las cajas del frente;
      en gris, los de atrás.}
    \end{column}
  \end{columns}
\end{frame}
% ============================================================
\begin{frame}{Dibujar un punto en el espacio}
  \begin{columns}[T]
    \begin{column}{0.52\textwidth}
      Para dibujar $P=(2,-1,2)$ caminamos desde el origen:

      \vspace{0.2cm}
      \begin{enumerate}
        \item $2$ unidades al frente (eje $x$)
        \item $1$ unidad a la izquierda (pues $y=-1$)
        \item $2$ unidades hacia arriba (eje $z$)
      \end{enumerate}

      \vspace{0.3cm}
      \destaca{\small Conviene \textbf{puntear la trayectoria} en el
      dibujo. Sin ella, la perspectiva engaña y el punto parece estar
      en otro lugar. Con ella, cualquiera puede reconstruir dónde
      vive.}
    \end{column}
    \begin{column}{0.44\textwidth}
      \begin{center}
      \begin{tikzpicture}[scale=1.0,>=Stealth]
        \draw[->,thick] (0,0) -- (-1.9,-1.52) node[below left]{$x$};
        \draw[->,thick] (0,0) -- (2.4,0) node[right]{$y$};
        \draw[gray!60] (0,0) -- (-1.7,0);
        \draw[->,thick] (0,0) -- (0,2.4) node[above]{$z$};
        \coordinate (A) at (-1.1,-0.88);
        \coordinate (B) at (-2.1,-0.88);
        \coordinate (P) at (-2.1,1.12);
        \draw[dashed,terracota,thick] (0,0) -- (A);
        \node[below right,font=\scriptsize] at (A) {$2$};
        \draw[dashed,terracota,thick] (A) -- (B);
        \draw[dashed,terracota,thick] (B) -- (P)
          node[midway,left,font=\scriptsize]{$2$};
        \draw[dashed,dorado,thick] (B) -- (-1,0);
        \node[above,font=\scriptsize] at (-1,0) {$-1$};
        \fill[dorado] (B) circle(1.8pt)
          node[below,font=\scriptsize]{$(2,-1,0)$};
        \fill[turquesaGA] (P) circle(2.6pt)
          node[above right,font=\small]{$P=(2,-1,2)$};
        \node[below left,font=\small] at (0,0) {$O$};
      \end{tikzpicture}
      \end{center}
    \end{column}
  \end{columns}
\end{frame}
% ============================================================
\begin{frame}{¿Dónde vive cada punto?}
  Como en $\mathbb{R}^2$, cerramos preguntando y dibujando cada
  respuesta:

  \vspace{0.2cm}
  \begin{enumerate}
    \item ¿En qué octante vive $(1,2,3)$? ¿Y $(2,-1,3)$?
    \item ¿Dónde vive el punto $(0,0,5)$?
    \item ¿Cuál es el lugar geométrico de los puntos con $z=0$?
    \item ¿Cuál es el lugar geométrico de los puntos con $y=3$?
  \end{enumerate}

  \vspace{0.4cm}
  \destaca{\small \textit{Respuestas.} 1.~Octante I; octante IV
  ($x>0$, $y<0$, $z>0$).\quad 2.~En el eje $z$; no está en ningún
  octante.\quad 3.~El \textbf{plano coordenado $xy$}: todos los puntos
  de altura cero, y solamente ellos.\quad 4.~El \textbf{plano paralelo
  al plano $xz$} que pasa por $(0,3,0)$. Las ecuaciones de estos
  planos las trabajaremos en la siguiente sesión.}
\end{frame}
% ============================================================
\section{Planos coordenados y distancia en el espacio}
% ---------- hoja de sesión ----------
\begin{frame}[plain]
  \hojaSesion{Planos coordenados y distancia en $\mathbb{R}^3$}{%
  Las ecuaciones de los planos coordenados y de los planos paralelos
  a ellos. La extensión de la fórmula de la distancia  al espacio (el mismo Pitágoras, aplicado dos veces).}
\end{frame}
\begin{frame}{Ecuaciones de los planos coordenados}
  El mismo razonamiento que en $\mathbb{R}^2$, donde el eje $X$ es la
  recta $y=0$:

  \vspace{0.2cm}
  \begin{center}
  {\small
  \begin{tabular}{@{}lcl@{}}
    \toprule
    \textbf{Plano coordenado} & \textbf{Ecuación} & \textbf{Descripción} \\
    \midrule
    plano $xy$ & $z=0$ & los puntos de altura cero \\
    plano $xz$ & $y=0$ & contiene a los ejes $x$ y $z$ \\
    plano $yz$ & $x=0$ & contiene a los ejes $y$ y $z$ \\
    \bottomrule
  \end{tabular}}
  \end{center}

  \vspace{0.3cm}
  \destaca{En $\mathbb{R}^2$, la ecuación $x=0$ deja libre una
  variable y describe una \textbf{recta} (el eje $Y$).
  
  \vspace{0.15cm}
  En $\mathbb{R}^3$, \textit{la misma ecuación} deja libres dos variables y
  describe un \textbf{plano}. 
  
  \vspace{0.15cm}
  Notemos que \textbf{la misma condición describe objetos
  distintos según el espacio donde viven los puntos}; esta observación
  reaparecerá varias veces en el curso.}
\end{frame}
% ============================================================
\begin{frame}{Planos paralelos a los coordenados}
  Como en el plano, donde $x=3$ es la recta vertical que pasa por
  $(3,0)$:

  \vspace{0.2cm}
  \begin{itemize}
    \item Los puntos con $y=3$ forman el \textbf{plano paralelo al
      plano $xz$} que pasa por $(0,3,0)$.
    \item Los puntos con $z=-2$ forman el plano paralelo al plano
      $xy$, dos unidades abajo.
    \item La variable que \textbf{no aparece} en la ecuación es la
      que queda libre; aquí quedan libres dos.
  \end{itemize}

  \vspace{0.3cm}
  \destaca{¿Y si imponemos \textbf{dos} condiciones a la vez? Los
  puntos con $y=3$ y $z=2$ forman una \textbf{recta} paralela al eje
  $x$ (la intersección de dos planos). Cada ecuación recorta un grado
  de libertad: en el espacio empezamos con tres, una condición deja
  un plano y dos condiciones dejan una recta.}
\end{frame}
% ============================================================
\begin{frame}{Distancia en el espacio desde un punto concreto hasta el origen}
  Tomamos $P=(2,6,3)$, en el primer octante, y bajamos la vertical
  desde $P$ hasta el piso. El pie es $P'=(2,6,0)$.

  \vspace{0.15cm}
  \begin{columns}[T]
    \begin{column}{0.40\textwidth}
      \begin{center}
      \begin{tikzpicture}[x={(-0.55cm,-0.44cm)},y={(1cm,0cm)},
                          z={(0cm,1cm)},scale=0.5,>=Stealth]
        % triangulos sombreados
        \fill[dorado!22] (0,0,0) -- (2,0,0) -- (2,6,0) -- cycle;
        \fill[turquesaGA!16] (0,0,0) -- (2,6,0) -- (2,6,3) -- cycle;
        % caja punteada
        \draw[gray!45,dashed] (2,6,0) -- (0,6,0) -- (0,0,0);
        \draw[gray!45,dashed] (0,0,3) -- (2,0,3) -- (2,6,3) -- (0,6,3) -- cycle;
        \draw[gray!45,dashed] (0,0,0) -- (0,0,3);
        \draw[gray!45,dashed] (2,0,0) -- (2,0,3);
        \draw[gray!45,dashed] (0,6,0) -- (0,6,3);
        % ejes
        \draw[->,gray!65] (0,0,0) -- (3.2,0,0) node[below,font=\scriptsize]{$x$};
        \draw[->,gray!65] (0,0,0) -- (0,7.6,0) node[right,font=\scriptsize]{$y$};
        \draw[->,gray!65] (0,0,0) -- (0,0,4.6) node[above,font=\scriptsize]{$z$};
        % triangulo del piso
        \draw[dorado,thick] (0,0,0) -- (2,0,0)
          node[midway,above left,font=\scriptsize]{$2$};
        \draw[dorado,thick] (2,0,0) -- (2,6,0)
          node[midway,below,font=\scriptsize]{$6$};
        \draw[dorado,very thick] (0,0,0) -- (2,6,0)
          node[pos=0.45,above,sloped,font=\scriptsize,yshift=1pt]{$\sqrt{40}$};
        \draw[dorado] (1.5,0,0) -- (1.5,0.5,0) -- (2,0.5,0);
        % triangulo vertical
        \draw[turquesaGA,thick] (2,6,0) -- (2,6,3)
          node[midway,right,font=\scriptsize]{$3$};
        \draw[turquesaGA,very thick] (0,0,0) -- (2,6,3)
          node[pos=0.6,above left,font=\scriptsize]{$7$};
        \draw[turquesaGA] (1.84,5.52,0) -- (1.84,5.52,0.5) -- (2,6,0.5);
        % puntos
        \fill[negroP] (0,0,0) circle(2pt) node[above left,font=\small]{$O$};
        \fill[turquesaGA] (2,6,3) circle(2.2pt) node[above,font=\small]{$P$};
        \fill[gray] (2,6,0) circle(1.8pt) node[below right,font=\scriptsize]{$P'$};
      \end{tikzpicture}
      \end{center}
    \end{column}
    \begin{column}{0.56\textwidth}
      \begin{tabular}{@{}c@{\hspace{8pt}}l@{}}
        \begin{tikzpicture}[scale=0.9]
          \fill[dorado!12] (0,0) -- (1.8,0) -- (1.8,0.6) -- cycle;
          \draw[dorado,thick] (0,0) -- (1.8,0)
            node[midway,below,font=\scriptsize]{$6$};
          \draw[dorado,thick] (1.8,0) -- (1.8,0.6)
            node[midway,right,font=\scriptsize]{$2$};
          \draw[dorado,very thick] (0,0) -- (1.8,0.6)
            node[pos=0.5,above,sloped,font=\scriptsize,yshift=1.5pt]{$\sqrt{40}$};
          \draw[dorado] (1.66,0) -- (1.66,0.14) -- (1.8,0.14);
        \end{tikzpicture}
        &
        \begin{minipage}{0.58\textwidth}
          \raggedright
          {\footnotesize\textcolor{dorado}{\textbf{Triángulo en el piso}}}\\[3pt]
          {\small $\overline{OP'}=\sqrt{2^2+6^2}=\sqrt{40}$}
        \end{minipage}
        \\[14pt]
        \begin{tikzpicture}[scale=0.9]
          \fill[turquesaGA!12] (0,0) -- (1.9,0) -- (1.9,0.9) -- cycle;
          \draw[turquesaGA,thick] (0,0) -- (1.9,0)
            node[midway,below,font=\scriptsize]{$\sqrt{40}$};
          \draw[turquesaGA,thick] (1.9,0) -- (1.9,0.9)
            node[midway,right,font=\scriptsize]{$3$};
          \draw[turquesaGA,very thick] (0,0) -- (1.9,0.9)
            node[pos=0.45,above,sloped,font=\scriptsize]{$7$};
          \draw[turquesaGA] (1.76,0) -- (1.76,0.14) -- (1.9,0.14);
        \end{tikzpicture}
        &
        \begin{minipage}{0.58\textwidth}
          \raggedright
          {\footnotesize\textcolor{turquesaGA}{\textbf{Triángulo vertical}}}\\[3pt]
          {\small $d(O,P)=\sqrt{\bigl(\sqrt{40}\bigr)^2+3^2}$}\\[2pt]
          {\small \phantom{$d(O,P)$}$=\sqrt{49}=7$}
        \end{minipage}
      \end{tabular}
    \end{column}
  \end{columns}

  \vspace{0.05cm}
  \destaca{\small El primer triángulo vive en el plano $xy$ y el
  segundo en el plano vertical que contiene a $O$, $P'$ y $P$. En
  ambos aplicamos Pitágoras, y el segundo toma como cateto la
  hipotenusa del primero.}
\end{frame}
% ============================================================
\begin{frame}{Distancia en el espacio desde un punto cualquiera hasta el origen}
  Repetimos el argumento con $P=(x_1,y_1,z_1)$ arbitrario. La
  proyección de $P$ sobre el plano $xy$ es $P'=(x_1,y_1,0)$.

  \vspace{0.15cm}
  \begin{columns}[T]
    \begin{column}{0.40\textwidth}
      \begin{center}
      \begin{tikzpicture}[x={(-0.55cm,-0.44cm)},y={(1cm,0cm)},
                          z={(0cm,1cm)},scale=0.5,>=Stealth]
        \fill[dorado!22] (0,0,0) -- (2,0,0) -- (2,6,0) -- cycle;
        \fill[turquesaGA!16] (0,0,0) -- (2,6,0) -- (2,6,3) -- cycle;
        % caja punteada
        \draw[gray!45,dashed] (2,6,0) -- (0,6,0) -- (0,0,0);
        \draw[gray!45,dashed] (0,0,3) -- (2,0,3) -- (2,6,3) -- (0,6,3) -- cycle;
        \draw[gray!45,dashed] (0,0,0) -- (0,0,3);
        \draw[gray!45,dashed] (2,0,0) -- (2,0,3);
        \draw[gray!45,dashed] (0,6,0) -- (0,6,3);
        % ejes
        \draw[->,gray!65] (0,0,0) -- (3.2,0,0) node[below,font=\scriptsize]{$x$};
        \draw[->,gray!65] (0,0,0) -- (0,7.6,0) node[right,font=\scriptsize]{$y$};
        \draw[->,gray!65] (0,0,0) -- (0,0,4.6) node[above,font=\scriptsize]{$z$};
        \draw[dorado,thick] (0,0,0) -- (2,0,0)
          node[midway,above left,font=\scriptsize]{$x_1$};
        \draw[dorado,thick] (2,0,0) -- (2,6,0)
          node[midway,below,font=\scriptsize]{$y_1$};
        \draw[dorado,very thick] (0,0,0) -- (2,6,0)
          node[pos=0.45,above,sloped,font=\scriptsize,yshift=1pt]{$\overline{OP'}$};
        \draw[dorado] (1.5,0,0) -- (1.5,0.5,0) -- (2,0.5,0);
        \draw[turquesaGA,thick] (2,6,0) -- (2,6,3)
          node[midway,right,font=\scriptsize]{$z_1$};
        \draw[turquesaGA,very thick] (0,0,0) -- (2,6,3)
          node[pos=0.6,above left,font=\scriptsize]{$d(O,P)$};
        \draw[turquesaGA] (1.84,5.52,0) -- (1.84,5.52,0.5) -- (2,6,0.5);
        \fill[negroP] (0,0,0) circle(2pt) node[above left,font=\small]{$O$};
        \fill[turquesaGA] (2,6,3) circle(2.2pt) node[above,font=\small]{$P$};
        \fill[gray] (2,6,0) circle(1.8pt) node[below right,font=\scriptsize]{$P'$};
      \end{tikzpicture}
      \end{center}
    \end{column}
    \begin{column}{0.56\textwidth}
      \begin{tabular}{@{}c@{\hspace{8pt}}l@{}}
        \begin{tikzpicture}[scale=0.9]
          \fill[dorado!12] (0,0) -- (1.8,0) -- (1.8,0.6) -- cycle;
          \draw[dorado,thick] (0,0) -- (1.8,0)
            node[midway,below,font=\scriptsize]{$|y_1|$};
          \draw[dorado,thick] (1.8,0) -- (1.8,0.6)
            node[midway,right,font=\scriptsize]{$|x_1|$};
          \draw[dorado,very thick] (0,0) -- (1.8,0.6)
            node[pos=0.5,above,sloped,font=\scriptsize,yshift=1.5pt]{$\overline{OP'}$};
          \draw[dorado] (1.66,0) -- (1.66,0.14) -- (1.8,0.14);
        \end{tikzpicture}
        &
        \begin{minipage}{0.58\textwidth}
          \raggedright
          {\footnotesize\textcolor{dorado}{\textbf{Triángulo en el piso}}}\\[3pt]
          {\small $\overline{OP'}=\sqrt{x_1^{\,2}+y_1^{\,2}}$}
        \end{minipage}
        \\[14pt]
        \begin{tikzpicture}[scale=0.9]
          \fill[turquesaGA!12] (0,0) -- (1.9,0) -- (1.9,0.9) -- cycle;
          \draw[turquesaGA,thick] (0,0) -- (1.9,0)
            node[midway,below,font=\scriptsize]{$\overline{OP'}$};
          \draw[turquesaGA,thick] (1.9,0) -- (1.9,0.9)
            node[midway,right,font=\scriptsize]{$|z_1|$};
          \draw[turquesaGA,very thick] (0,0) -- (1.9,0.9)
            node[pos=0.45,above,sloped,font=\scriptsize]{$d(O,P)$};
          \draw[turquesaGA] (1.76,0) -- (1.76,0.14) -- (1.9,0.14);
        \end{tikzpicture}
        &
        \begin{minipage}{0.58\textwidth}
          \raggedright
          {\footnotesize\textcolor{turquesaGA}{\textbf{Triángulo vertical}}}\\[3pt]
          {\small $d(O,P)=\sqrt{x_1^{\,2}+y_1^{\,2}+z_1^{\,2}}$}
        \end{minipage}
      \end{tabular}
    \end{column}
  \end{columns}

  \vspace{0.05cm}
  \destaca{\small Los catetos del piso miden $|x_1|$ y $|y_1|$, porque
  una longitud no es negativa. Los valores absolutos desaparecen al
  elevar al cuadrado, y por eso no aparecen en la fórmula.}
\end{frame}
% ============================================================
\begin{frame}{Distancia entre dos puntos del espacio}
  Damos el último paso con $P=(x_1,y_1,z_1)$ y $Q=(x_2,y_2,z_2)$.
  Trasladamos el triángulo igual que en el plano, tomando como piso
  la altura de $P$. El vértice auxiliar es $R=(x_2,y_2,z_1)$.

  \vspace{0.1cm}
  \begin{columns}[T]
    \begin{column}{0.42\textwidth}
      \begin{center}
      \begin{tikzpicture}[x={(-0.55cm,-0.44cm)},y={(1cm,0cm)},
                          z={(0cm,1cm)},scale=0.42,>=Stealth]
        % triangulos sombreados
        \fill[dorado!22] (1,4,3) -- (4,4,3) -- (4,8,3) -- cycle;
        \fill[turquesaGA!16] (1,4,3) -- (4,8,3) -- (4,8,6) -- cycle;
        % aristas de la caja, punteadas
        \draw[gray!50,dashed] (1,4,6) -- (4,4,6) -- (4,8,6) -- (1,8,6) -- cycle;
        \draw[gray!50,dashed] (1,4,3) -- (1,4,6);
        \draw[gray!50,dashed] (4,4,3) -- (4,4,6);
        \draw[gray!50,dashed] (1,8,3) -- (1,8,6);
        \draw[gray!50,dashed] (1,4,3) -- (1,8,3) -- (4,8,3);
        % ejes
        \draw[->,gray!55] (0,0,0) -- (3.2,0,0) node[below,font=\scriptsize]{$x$};
        \draw[->,gray!55] (0,0,0) -- (0,9.4,0) node[right,font=\scriptsize]{$y$};
        \draw[->,gray!55] (0,0,0) -- (0,0,7) node[above,font=\scriptsize]{$z$};
        % triangulo del piso
        \draw[dorado,thick] (1,4,3) -- (4,4,3)
          node[pos=0.62,above,sloped,font=\tiny,yshift=1pt]{$|x_2-x_1|$};
        \draw[dorado,thick] (4,4,3) -- (4,8,3)
          node[midway,below,font=\tiny]{$|y_2-y_1|$};
        \draw[dorado,very thick] (1,4,3) -- (4,8,3);
        \draw[dorado] (3.6,4,3) -- (3.6,4.5,3) -- (4,4.5,3);
        % triangulo vertical
        \draw[turquesaGA,thick] (4,8,3) -- (4,8,6)
          node[midway,below,sloped,font=\tiny,yshift=-2pt]{$|z_2-z_1|$};
        \draw[turquesaGA,very thick] (1,4,3) -- (4,8,6)
          node[pos=0.55,above left,font=\scriptsize]{$d(P,Q)$};
        \draw[turquesaGA] (3.83,7.5,3) -- (3.83,7.5,3.5) -- (4,8,3.5);
        % puntos
        \fill[turquesaGA] (1,4,3) circle(2.2pt) node[above right,font=\small,xshift=-2pt]{$P$};
        \fill[turquesaGA] (4,8,6) circle(2.2pt) node[above,font=\small]{$Q$};
        \fill[gray] (4,8,3) circle(1.8pt) node[below right,font=\scriptsize]{$R$};
      \end{tikzpicture}
      \end{center}
    \end{column}
    \begin{column}{0.54\textwidth}
      \begin{center}
      \begin{tabular}{@{}c@{\hspace{16pt}}c@{}}
        \begin{tikzpicture}[scale=0.72]
          \fill[dorado!12] (0,0) -- (1.7,0) -- (1.7,0.85) -- cycle;
          \draw[dorado,thick] (0,0) -- (1.7,0)
            node[midway,below,font=\tiny]{$|y_2-y_1|$};
          \draw[dorado,thick] (1.7,0) -- (1.7,0.85)
            node[midway,right,font=\tiny]{$|x_2-x_1|$};
          \draw[dorado,very thick] (0,0) -- (1.7,0.85)
            node[midway,above,sloped,font=\tiny]{$\overline{PR}$};
          \draw[dorado] (1.55,0) -- (1.55,0.15) -- (1.7,0.15);
        \end{tikzpicture}
        &
        \begin{tikzpicture}[scale=0.72]
          \fill[turquesaGA!12] (0,0) -- (1.7,0) -- (1.7,1.3) -- cycle;
          \draw[turquesaGA,thick] (0,0) -- (1.7,0)
            node[midway,below,font=\tiny]{$\overline{PR}$};
          \draw[turquesaGA,thick] (1.7,0) -- (1.7,1.3)
            node[midway,right,font=\tiny]{$|z_2-z_1|$};
          \draw[turquesaGA,very thick] (0,0) -- (1.7,1.3)
            node[midway,above,sloped,font=\tiny]{$d(P,Q)$};
          \draw[turquesaGA] (1.55,0) -- (1.55,0.15) -- (1.7,0.15);
        \end{tikzpicture}
        \\[3pt]
        {\footnotesize\textcolor{dorado}{\textbf{En el piso}}} &
        {\footnotesize\textcolor{turquesaGA}{\textbf{Vertical}}}
      \end{tabular}
      \end{center}

      \vspace{0.1cm}
      {\footnotesize
      \[ \overline{PR}=\sqrt{(x_2-x_1)^2+(y_2-y_1)^2} \]

      \vspace{-0.2cm}
      \[ d(P,Q)=\sqrt{(x_2-x_1)^2+(y_2-y_1)^2+(z_2-z_1)^2} \]}
    \end{column}
  \end{columns}

  \vspace{0.05cm}
  \destaca{\small \textbf{Ejemplo 5.}
  $d\bigl((2,-1,3),\,(0,1,2)\bigr)=\sqrt{4+4+1}=\sqrt{9}=3$. Si $P$ es
  el origen recuperamos la fórmula anterior, y si $z_1=z_2$ recuperamos
  la del plano.}
\end{frame}
% ============================================================
\begin{frame}{Las propiedades también en el espacio}
  Las tres propiedades que verificamos en el plano se conservan en
  $\mathbb{R}^3$, y se verifican igual, porque la fórmula sólo agrega
  un sumando.

  \vspace{0.25cm}
  \begin{itemize}
    \item[1.] $d(P,Q)\ \ge\ 0$, \ y \ $d(P,Q)=0 \iff P=Q$. Una raíz de
      suma de cuadrados no es negativa, y vale $0$ únicamente cuando
      los \textbf{tres} cuadrados son $0$, es decir cuando $x_1=x_2$,
      $y_1=y_2$ y $z_1=z_2$.
    \item[2.] $d(P,Q)=d(Q,P)$. Elevar al cuadrado elimina el signo,
      también en el tercer sumando.
    \item[3.] $d(P,Q)\ \le\ d(P,R)+d(R,Q)$. Sigue \textbf{enunciada y
      no demostrada}, ahora también en el espacio.
  \end{itemize}

  \vspace{0.35cm}
  \destaca{Esta es la ventaja de haber construido $\mathbb{R}^3$ en
  paralelo con $\mathbb{R}^2$. Casi nada se demuestra otra vez, basta
  revisar que el argumento anterior siga funcionando con una
  coordenada más. La desigualdad del triángulo la demostraremos más adelante, y una sola demostración servirá para los dos casos.}
\end{frame}
% ============================================================
\begin{frame}{Cierre de la semana}
  \textbf{Lo que vimos}
  {\small
  \begin{itemize}
    \item La correspondencia figura $\leftrightarrow$ ecuación y el
      concepto de lugar geométrico %(\textit{todos} y \textit{solamente})

    \item El plano cartesiano, los signos y la distancia como teorema
      de Pitágoras %(una propiedad quedó pendiente de demostrar)
    \item El espacio $\mathbb{R}^3$ en paralelo con $\mathbb{R}^2$:
      octantes, planos coordenados y sus ecuaciones
    \item La distancia en $\mathbb{R}^3$, con su fórmula y los
      primeros cálculos %(la deducción con Pitágoras la vemos en el      video y la retomamos al abrir la próxima sesión)
  \end{itemize}}

  \vspace{0.05cm}
  \textbf{Para la próxima semana}
  {\small
  \begin{itemize}
    \item Ver el video \textit{Distancia entre dos puntos en el
      espacio tridimensional} (TURBO CLASES), donde se deduce la
      fórmula aplicando dos veces el teorema de Pitágoras\\
      {\footnotesize\url{https://youtu.be/iqiO7iMouQE}}
    \item %Lectura: Ramírez-Galarza, sección 1.
    Visitar GeoGebra 3D y
      mover la perspectiva
    \item Cuestionario de práctica de la semana 1 en Moodle
      (intentos ilimitados; no forma parte de la calificación, pero son el pase para el parcial)
  \end{itemize}}

%  Respuesta de la pregunta oculta: si $xy>0$, las coordenadas tienen
%  el mismo signo; si además $x+y<0$, ambas son negativas, así que el
%  punto está en el cuadrante III. La segunda pregunta abre la semana 2
%  (los puntos que equidistan de dos puntos dados forman la mediatriz).
%
%  \vspace{0.2cm}
%  \destaca{\small \textbf{Pregunta para la próxima sesión:} si $xy>0$
%  y $x+y<0$, ¿en qué cuadrante se encuentra el punto $(x,y)$?
%  Y una más de fondo: ¿qué lugar geométrico resulta si la condición
%  involucra distancias a \textbf{dos} puntos? Con esa pregunta abrimos
%  la semana 2.}
\end{frame}
% ============================================================
%  Página de cierre del archivo.
%  Ajustar la URL cuando el libro quede publicado en GitHub Pages.
\begin{frame}[plain]
  \centering

  \vspace{1.0cm}
  \begin{tikzpicture}[scale=0.55,>=Stealth]
    \draw[very thin,color=gray!30] (-2.6,-3.1) grid (4.1,2.6);
    \draw[->,thick] (-2.6,0) -- (4.3,0) node[below,font=\small]{$x$};
    \draw[->,thick] (0,-3.1) -- (0,2.8) node[left,font=\small]{$y$};
    \draw[turquesaGA,very thick] (3,-3.1) -- (3,2.6);
    \draw[terracota,very thick] (-2.6,-2) -- (4.1,-2);
    \fill[dorado] (3,-2) circle(3pt);
  \end{tikzpicture}

  \vspace{0.5cm}

  {\Large\color{turquesaGA} ¿Preguntas?}\\[10pt]
  
  \vspace{0.7cm}

  {\tiny\color{grisT}
  Texto base: Ramírez Galarza, \textit{Geometría analítica. Una
  introducción a la geometría}, UNAM. ---
  Ejercicios: Lehmann, \textit{Geometría analítica}, Limusa.}
\end{frame}
\end{document}
